% from ega4-21b-fr.tex
\oldpage[IV-4]{286}
Pour le 1\textsuperscript{o}, il suffit de noter que pour tout ouvert $U$ de $X$ et tout diviseur $D\in\operatorname{Div}(U)$, les points maximaux du support $S$ de $D$ sont tels que $\operatorname{prof}(\sh{O}_{X,x})=1$ \sref{IV.21.1.9}, donc \emph{a fortiori} appartiennent à $X^{(1)}\cap U$;
puisque $\dim(X)\leq1$, l'ensemble de ces points est égal à $S$, et $S$ est donc discret, l'ensemble des composantes irréductibles de $S$ étant localement fini.
\end{proof}

\begin{corollary}[21.9.5]
\label{IV.21.9.5}
Soit $X$ un préschéma localement noethérien de dimension $\leq1$.
Pour toute partie discrète $E\subset X^{(1)}$ et toute famille $(D(x))_{x\in E}$ avec $D(x)\in\operatorname{Div}(\sh{O}_{X,x})$, il existe un diviseur $D$ sur $X$ et un seul tel que le support de $D$ soit contenu dans $E$ et que $D_x=D(x)$ pour tout $x\in E$.
\end{corollary}

\begin{proof}
Cela résulte de \sref{IV.21.9.4} et de \sref{IV.21.9.3}, (ii).
\end{proof}

\begin{corollary}[21.9.6]
\label{IV.21.9.6}
Soit $X$ un préschéma localement noethérien de dimension $\leq1$;
tout diviseur $D$ sur $X$ est de la forme $D'-D''$, où $D'$, $D''$ sont deux diviseurs $\geq0$ de supports contenus dans celui de $D$.
\end{corollary}

\begin{proof}
En vertu de \sref{IV.21.9.5}, on est aussitôt ramené au cas où $X=\Spec(A)$, où $A$ est un anneau local noethérien;
il suffit alors d'observer que $M(X)$ est l'anneau total des fractions de $A$, et qu'une section de $\sh{M}^*_X$ au-dessus de $X$ s'écrit par définition comme quotient $b/a$ de deux éléments réguliers de $A$.
\end{proof}

\begin{corollary}[21.9.7]
\label{IV.21.9.7}
Soient $X$ un préschéma localement noethérien de dimension $\leq1$, sans point isolé, et $U$ un ouvert dense dans $X$.
Il existe alors un diviseur $D\geq0$ sur $X$, de support contenu dans $U$ et rencontrant chacune des composantes irréductibles de $X$.
\end{corollary}

\begin{proof}
On peut supposer que $U$ est réunion d'ensembles ouverts disjoints non vides $U_\alpha$, dont chacun est contenu dans une seule composante irréductible de $X$;
il suffit alors de prendre dans chaque $U_\alpha$ un point $x_\alpha$ fermé dans $X$ (il en existe puisque l'unique point générique de $U_\alpha$ ne peut être isolé par hypothèse), et d'appliquer \sref{IV.21.9.5} à l'ensemble discret des $x_\alpha$.
\end{proof}

L'intérêt de ce corollaire est qu'il permettra de prouver qu'une courbe algébrique séparée $X$ sur un corps $k$ est quasi-projective, le diviseur $D\geq0$ défini dans \sref{IV.21.9.7} étant alors nécessairement \emph{ample} en vertu du théorème de Riemann--Roch pour les courbes (chap.~V).

\begin{env}[21.9.8]
\label{IV.21.9.8}
Pour les préschémas localement noethériens $X$ de dimension $\leq1$ la proposition \sref{IV.21.9.4} ramène la détermination de $\sh{D}iv_X$ au cas où $X=\Spec(A)$, $A$ étant un anneau local noethérien de dimension $1$.

Lorsque $A$ est un anneau local régulier de dimension $1$ (autrement dit, un anneau de valuation discrète), le groupe $\operatorname{Div}(A)$ s'identifie canoniquement à $\bb{Z}$ \sref{IV.21.6.8};
par suite, en vertu de \sref{IV.21.9.2}:
\end{env}

\begin{proposition}[21.9.9]
\label{IV.21.9.9}
Soit $X$ un préschéma localement noethérien régulier de dimension $\leq1$.
Alors le faisceau $\sh{D}iv_X$ est canoniquement isomorphe à $\bigoplus_{x\in X^{(1)}}(i_x)_*(\bb{Z}(x))$, où $\bb{Z}(x)$ est le groupe additif $\bb{Z}$ considéré comme faisceau de groupes sur l'espace $\{x\}$.
\end{proposition}

\begin{env}[21.9.10]
\label{IV.21.9.10}
Supposons seulement que l'anneau local noethérien $A$ de dimension $1$ soit \emph{réduit};
alors, si $A'$ est le \emph{normalisé} de $A$ (fermeture intégrale de $A$ dans son anneau total de fractions), $A'$ est noethérien en vertu du théorème de Krull--Akizuki (Bourbaki, \emph{Alg.\ comm.}, chap.~VII, \textsection~2, n\textsuperscript{o}~5, prop.~5), et on a vu dans \sref{IV.21.8.6} comment
